\section{Empirical Evidence and Falsifiable Predictions}

\subsection{What is already established}

ARC builds on components that should not be presented as novel. Attachment anxiety and avoidance are established dimensional descriptions of adult attachment, and meta-analytic evidence links both dimensions to romantic relationship quality \parencite{fraley2000irt,mikulincer2016attachment,li2012meta}. Positive secure-base representations have an independent theoretical and measurement tradition \parencite{waters2006workingmodels,waters2019securebase,waters2025securebase}. Control-system interpretations of attachment predate ARC \parencite{bowlby1982attachment,kobak1993control,petters2019control}, and computational models already implement targets, state updating, feedback, and recovery-like dynamics \parencite{stevens2009dynamic,gagliardi2022control,olar2026relationship}. Adult relationship theories also provide established accounts of risk regulation, coregulation, and partner buffering of attachment insecurity \parencite{murray2006risk,sbarra2008coregulation,simpson2014buffering,overall2022buffering}.

The empirical contribution of ARC therefore depends on joining these strands in a single temporal measurement architecture and testing whether the proposed dynamic variables add information beyond simpler alternatives.

\subsection{Relational discrepancy has empirical precedent, but H1 remains under-tested}

The distinction between experienced and desired closeness has direct empirical precedent. Longitudinal work using actual and ideal closeness showed that smaller discrepancies were associated with better relational well-being and mental health, and that moving closer to one's ideal over time accompanied improved outcomes \parencite{frost2013closeness}. Dyadic response-surface analyses likewise showed that closeness discrepancy is consequential for relationship quality, with especially unfavorable outcomes when experienced closeness is lower than desired \parencite{pusch2022closeness}. Population-based evidence further indicates that wanting more closeness is more strongly associated with poorer relationship characteristics than experiencing more closeness than desired \parencite{pietras2024closeness}.

These findings support the construct validity of a perceived--desired discrepancy, but they do not establish the ARC-specific causal-temporal claim. Most discrepancy studies predict relationship quality, well-being, or longer-term change rather than the next regulatory action within the same individual. H1 is therefore not the claim that closeness discrepancies exist or matter. It is the narrower claim that momentary discrepancy contains incremental information about subsequent regulation after current state, context, and behavioral inertia have been accounted for.

\subsection{Why the difference score must itself compete}

A simple score
\[
e_t=r^\ast_t-\hat r_t
\]
imposes a substantive constraint: desired state and perceived state are assumed to enter the action model with equal magnitude and opposite sign. That constraint should be tested rather than hidden inside data preprocessing.

The H1 study therefore compares a sequence of increasingly informative observation models.

\begin{table}[htbp]
\caption{Competing Observation Models for the H1 Discrepancy Claim}
\label{tab:h1-competing-models}
\centering
\small
\setlength{\tabcolsep}{3pt}
\renewcommand{\arraystretch}{0.95}
\begin{threeparttable}
\begin{tabular}{p{0.09\linewidth}p{0.30\linewidth}p{0.24\linewidth}p{0.24\linewidth}}
\toprule
Model & Additional relational predictors & Question answered & Interpretation \\
\midrule
M0 & None beyond context and prior action & Is generic context plus behavioral inertia sufficient? & Baseline model \\
M1 & Perceived state \(\hat r_t\) & Does current relational perception alone predict action? & Descriptive baseline for relational-state information \\
M2 & \(\hat r_t\) and \(r^\ast_t\) separately & Does desired state contribute beyond perceived state? & Flexible component model; \(\beta_{r^\ast}\) is H1-P1 \\
M3 & Discrepancy \(e_t=r^\ast_t-\hat r_t\) & Does signed discrepancy predict subsequent action? & \(\beta_e\) is H1-P2; equality is checked separately in M2 \\
M4 & Positive and negative discrepancy components & Are closeness frustration and closeness surfeit dynamically asymmetric? & Piecewise discrepancy extension \\
\bottomrule
\end{tabular}
\begin{tablenotes}[flushleft]
\footnotesize
\item Note. All models retain the same preregistered context, elapsed-time, and autoregressive terms. H1-P1 and H1-P2 are direct coefficient tests with Holm family-wise control. Structural contrasts are evaluated separately. Days 12--14 remain untouched for secondary predictive RMSE validation.
\end{tablenotes}
\end{threeparttable}
\end{table}


The discrepancy representation is strengthened if the constrained \(e_t\) model predicts future action at least as well as the less parsimonious unconstrained model containing \(r^\ast_t\) and \(\hat r_t\) separately. It is weakened if perceived state alone performs equally well, because the target adds no useful information, or if the two raw components have materially different effects that are obscured by subtraction.

Existing evidence also makes directional asymmetry plausible. Wanting more closeness than currently experienced may not be dynamically equivalent to wanting less closeness than currently experienced \parencite{pusch2022closeness,pietras2024closeness}. ARC therefore treats positive and negative discrepancy components as a mandatory sensitivity analysis rather than assuming one globally linear slope.

\subsection{Perceived partner responsiveness supports the feedback pathway}

The feedback component also has strong adjacent evidence. In a two-week daily diary study, within-person increases in perceived partner responsiveness covaried with greater emotional expression to the partner; an experimental study in the same program showed that induced responsiveness increased expression of anxiety \parencite{ruan2020responsiveness}. A yearlong intensive longitudinal study further showed that perceived partner responsiveness participated in within- and between-person commitment dynamics and was associated with attachment-related individual differences \parencite{segal2016responsiveness}.

These studies make the ARC feedback pathway plausible, but they also highlight a measurement distinction that the framework must preserve. Perceived responsiveness is not identical to partner behavior. It may represent a psychologically important state estimate while also containing projection, attachment-related interpretation, and common-method variance. Stronger dyadic tests should therefore distinguish actor perception from partner report or independently coded partner behavior whenever possible.

\subsection{Attachment orientations plausibly shape both response magnitude and response channel}

The idea that \Anx\ and \Avo\ modify relational dynamics has direct within-person precedent and is consistent with interactionist accounts in which attachment effects depend on relational context and activation conditions \parencite{campbell2011interactionist}. Event-contingent work found that attachment anxiety strengthened, and attachment avoidance weakened, affective reactivity to agreeable interpersonal behavior, particularly with romantic partners \parencite{sadikaj2011dynamics}. Daily and intensive longitudinal studies likewise show attachment-related differences in relationship feedback, commitment processes, and responses to relational cues \parencite{segal2016responsiveness,beck2024daily,carnelley2025phubbing}. More recent dyadic experience-sampling work using dynamic structural equation modeling reports attachment-related moderation of emotional interdependence between partners \parencite{xie2025interdependence}.

This evidence argues against representing attachment only as a shift in one signed action coefficient. Anxious regulation can be expressed through overt proximity seeking, reassurance seeking, protest, or heightened interpersonal engagement, while avoidant regulation can appear as weaker overt responding, distancing, suppression, or internal disengagement. The same relational discrepancy can therefore be translated into different response channels.

ARC consequently treats H2 as a multivariate policy claim. The key test is whether \Anx\ and \Avo\ explain reproducible heterogeneity in channel-specific discrepancy-to-action mappings after context and action history are controlled. The strongest directional prediction concerns positive discrepancy, where a person wants more closeness than is currently experienced: higher \Anx\ should preferentially strengthen overt approach-oriented responses, whereas higher \Avo\ should reduce overt approach and increase distancing or internal disengagement. Effects for closeness surfeit remain less constrained and should be estimated rather than assumed.

This literature makes H2 plausible, but does not establish the ARC-specific \(e_t\rightarrow a_{t+1}\) moderation. H2 therefore remains contingent on H1 measurement first proving viable.

\subsection{Secure-base processes justify a candidate, not yet a required core construct}

Adult secure-base functioning can be measured behaviorally. In engaged couples, the Secure Base Scoring System captured secure-base use and support in observed problem-solving interactions and predicted relationship variables beyond a generic communication coding system \parencite{crowell2002securebase}. Daily dyadic work likewise found reciprocal associations between relationship-specific security and support seeking, provision, receipt, and felt support \parencite{davila2009securebase}.

There is also direct evidence that attachment is associated with recovery from relational stress. Cortisol trajectories after relationship conflict vary with attachment insecurity \parencite{powers2006cortisol,pietromonaco2013cortisol}, positive dyadic coping predicts faster cortisol recovery and interacts with attachment anxiety \parencite{meuwly2012cortisol}, and observed recovery from romantic conflict has developmental links to earlier attachment security and partner functioning \parencite{salvatore2011recovery}. Coregulatory and partner-buffering accounts independently place interpersonal regulation and recovery at the center of adult attachment processes \parencite{sbarra2008coregulation,simpson2014buffering,overall2022buffering}. Relationship experiences can also contribute to longer-term changes in attachment security, supporting a distinction between momentary regulation and slower representational change \parencite{arriaga2019security,bayraktaroglu2023positive}.

However, this evidence does not yet demonstrate the specific ARC H3 claim. Much of the recovery literature operationalizes security through low insecurity, attachment classifications, developmental history, or partner behavior. The direct incremental question remains open: does a positive secure-base representation \Sec\ predict recovery after \Anx\ and \Avo\ are already modeled?

H3 is therefore deliberately stronger than the existing evidence and serves as a promotion test rather than as an assumption required to define ARC. \Sec\ should be promoted to the person-level core only if a positive measure of secure-base representation improves a preregistered primary recovery outcome beyond low attachment anxiety and avoidance, adds genuinely out-of-sample predictive information, and the increment is independently replicated or supported by a second preregistered design. The hierarchy of tests is summarized in Table~\ref{tab:h2-h3-tests}.

\begin{table}[htbp]
\caption{Predefined Tests for Attachment-Dependent Policy and Secure-Base Recovery}
\label{tab:h2-h3-tests}
\centering
\scriptsize
\setlength{\tabcolsep}{2pt}
\renewcommand{\arraystretch}{0.85}
\begin{threeparttable}
\begin{tabular}{p{0.09\linewidth}p{0.26\linewidth}p{0.27\linewidth}p{0.26\linewidth}}
\toprule
Test & Model contrast & Target quantity & Decision rule \\
\midrule
H2a & Channel-specific action model without versus with \(\Anx \times e_t\) & Approach/proximity-seeking response to closeness deficit & Retain only if moderation is reproducible and improves held-out prediction \\
H2b & Channel-specific action model without versus with \(\Avo \times e_t\) & Overt approach, distancing, and internal disengagement & Retain channel-specific effects; do not force one signed-action interpretation \\
H2c & Actor-only versus actor--partner moderation & Partner attachment contribution to actor response & Treat as dyadic extension unless incremental fit and prediction justify it \\
H3a & Recovery baseline versus baseline plus \(\Anx,\Avo\) & Recovery explained by insecurity dimensions & Establishes comparison model, not positive security \\
H3b & H3a plus \(\Sec\) & Incremental value of positive secure-base representation & Promotion requires primary-outcome improvement beyond \(\Anx,\Avo\), out-of-sample value, and replication; target retreat alone is not recovery \\
H3c & H3b plus observed/perceived partner support & Whether \(\Sec\) remains informative after proximal support is modeled & Stronger evidence for independent core status; otherwise report conditional attenuation, not identified mediation \\
\bottomrule
\end{tabular}
\begin{tablenotes}[flushleft]
\footnotesize
\item Note. H2 is evaluated on separate action channels as well as any signed summary. H3 is a promotion test for a positive secure-base measure; secondary recovery outcomes cannot rescue failure on the preregistered primary endpoint.
\end{tablenotes}
\end{threeparttable}
\end{table}


\subsection{Core hypotheses and candidate-core promotion test}

\textbf{H1: Discrepancy--action coupling.} Within individuals, measured desired-minus-perceived relational discrepancy at time \(t\) predicts subsequent regulatory action at \(t+1\) beyond measured context, perceived state alone, and autoregressive action history.

\textbf{H2: Attachment-dependent policy.} Attachment anxiety and avoidance explain reproducible between-person differences in how discrepancy and context are translated into specific regulatory channels. Anxiety is expected to preferentially strengthen overt approach-oriented responses to closeness deficit, whereas avoidance is expected to reduce overt approach and increase distancing or internal disengagement. The prediction concerns moderation of measured temporal relations, not assignment of fixed controller types.

\textbf{H3: Secure-base promotion test.} Positive secure-base representation predicts more effective recovery and return to constructive activity beyond what is explained by low anxiety and avoidance and, in stronger tests, beyond proximal partner support. H3 is the preregistered criterion for promoting \Sec\ from candidate-core status to the required person-level core.

\subsection{Extension hypotheses}

\textbf{H4: Candidate moderator specificity.} Mentalized affectivity predicts or moderates externally scored relational-state inference quality, while emotion-regulation flexibility predicts or moderates context- and feedback-sensitive action adaptation. These effects must be incremental to adjacent constructs and are not required for the core ARC loop.

\textbf{H5: Dynamic organization and external disorganization.} Preregistered dynamic indicators such as poor recovery, context-residual inconsistency, persistent approach--avoidance conflict, or maladaptive switching improve held-out prediction of independently measured adult disorganization beyond \Anx, \Avo, their interaction, and simpler static summaries.

\subsection{Evidence hierarchy}

The five hypotheses do not have equal evidential status. H1 is the necessary first test because the feedback architecture loses its empirical purpose if measured discrepancy does not predict subsequent action. H2 tests whether established attachment dimensions modulate the policy. H3 is a candidate-core promotion test for positive security, not a requirement for the current minimal model. H4 and H5 are extensions whose failure should simplify the framework rather than trigger redefinition of the core.

\begin{table}[htbp]
\caption{Evidence Status and Decisive Tests for ARC Hypotheses}
\label{tab:hypothesis-evidence}
\centering
\scriptsize
\setlength{\tabcolsep}{2pt}
\renewcommand{\arraystretch}{0.85}
\begin{threeparttable}
\begin{tabular}{p{0.08\linewidth}p{0.27\linewidth}p{0.27\linewidth}p{0.26\linewidth}}
\toprule
Hypothesis & Adjacent empirical support & ARC-specific evidence still missing & Decisive result \\
\midrule
H1 & Closeness discrepancies predict relational outcomes \parencite{frost2013closeness,pusch2022closeness,pietras2024closeness} & Within-person lag from measured discrepancy to subsequent regulatory action & Target/state information improves held-out action prediction and the retained discrepancy representation survives comparison with simpler alternatives \\
H2 & Attachment moderates within-person reactivity, daily relational responses, and dyadic interdependence \parencite{sadikaj2011dynamics,segal2016responsiveness,xie2025interdependence,carnelley2025phubbing} & Direct moderation of channel-specific \(e_t\rightarrow a^{(k)}_{t+1}\) mappings & Reproducible cross-level moderation of approach, distance, or internal regulation after H1 measurement is established \\
H3 candidate core & Adult secure-base behavior is measurable; attachment and dyadic coping predict support and stress recovery \parencite{crowell2002securebase,davila2009securebase,powers2006cortisol,meuwly2012cortisol,salvatore2011recovery} & Incremental prediction of preregistered recovery outcomes by positive \(S\) beyond \(L/U\) and, in stronger tests, proximal support & Promote \(S\) only after preregistered primary-outcome improvement beyond insecurity, out-of-sample value, and replication \\
H4 & Mentalized affectivity and flexibility have construct-level support \parencite{greenberg2017mentalized,mosannenzadeh2024flexibility} & Pathway-specific incremental validity in repeated relational data & Moderator improves its assigned pathway beyond adjacent constructs \\
H5 & Adult disorganization is measurable but heterogeneous \parencite{pollard2023disorganized} & Incremental validity of dynamic trajectory markers & Dynamic markers predict an independent criterion beyond static attachment baselines \\
\bottomrule
\end{tabular}
\begin{tablenotes}[flushleft]
\footnotesize
\item Note. Adjacent support is not treated as confirmation of the ARC-specific temporal hypothesis.
\end{tablenotes}
\end{threeparttable}
\end{table}


\subsection{Comparison with prior models}

\begin{table}[htbp]
\caption{Selected Precedents and the Empirical Position of ARC}
\label{tab:model-comparison}
\centering
\scriptsize
\setlength{\tabcolsep}{2pt}
\renewcommand{\arraystretch}{0.85}
\begin{threeparttable}
\begin{tabular}{p{0.17\linewidth}p{0.23\linewidth}p{0.23\linewidth}p{0.25\linewidth}}
\toprule
Approach & Existing contribution & Human dynamic evidence & ARC empirical role \\
\midrule
Bowlby / Kobak & Goal-directed attachment control and emotion-regulation set-goals & Observational mother--teen work for Kobak et al. & Does not claim novelty; translates control assumptions into repeated adult measures \\
Stevens and Zhang & Continuous dynamical attachment system with recovery-like behavior & Simulation / parameter analysis & Psychological rather than neurochemical state measurement in adults \\
Petters & Explicit computational lineage of the attachment control system & Primarily theoretical / computational & Uses the lineage as prior theory, not as an ARC-specific mechanism \\
Gagliardi & Multidimensional set-goals, perception, action, feedback, and agent interaction & Computational implementation & Tests a comparable temporal logic with intensive human observations \\
Murray / Sbarra and Hazan & Risk regulation, coregulation, separation and recovery in adult relationships & Human relationship literature and empirical agenda & Integrates adult interpersonal regulation with explicit time-ordered measurement \\
Olar et al. & Adaptive network account of anxious--avoidant dyadic dynamics & Computational model & Provides a rival dynamic architecture for comparison \\
ARC & Parsimonious sequence from measured state and target to action, feedback, and recovery & Proposed intensive longitudinal and dyadic tests & Contribution is integration, measurement architecture, held-out prediction, and falsification \\
\bottomrule
\end{tabular}
\begin{tablenotes}[flushleft]
\footnotesize
\item Note. The table identifies differences in scope and evidential role. It is not a ranking of models.
\end{tablenotes}
\end{threeparttable}
\end{table}


The comparison emphasizes that ARC does not originate control-system attachment. Its proposed addition is a common measurement bridge from person-level attachment variables to repeated relational states, actions, partner feedback, and recovery in adult relationships.
